
% \section{Empirical Study - Understanding Control Flows in Hacked DeFi Protocols} 
% \label{sec:study}


\begin{figure*}[!h]
\centering
\includegraphics[width=0.8\textwidth]{figures/guard_new_new.pdf}
    \vspace{-2mm}
    \caption{Running example showing \textsc{CrossGuard}'s response to simple invocation 
    (left) and re-entrancy attack (right)}
    \Description{Two-panel running example. Left: a benign/simple invocation sequence. Right: a re-entrancy attack sequence. The figure illustrates how \textsc{CrossGuard} observes function-level control flow and reacts differently in benign versus attack cases.}
    \label{fig:example}
\end{figure*}


To ground our framework, we hypothesize that \emph{function-level} 
control flows are sufficient to distinguish malicious from 
benign transactions, while remaining practical overheads
for runtime monitoring. 
To evaluate this hypothesis, we analyze each 
victim protocol's 
transaction history prior to the incident. 
Because protocols typically 
operate for some time 
before being exploited and attract many 
interacting actors, benign behavior can 
evolve over time, diversifying and introducing 
previously unseen control flows. 
This motivates the following research question:



\noindent \textbf{RQ1: 
How do \revise{(function level)} control flows in hack transactions differ 
from those in other 
(benign) transactions 
prior to a hack?}

To answer RQ1, we conducted a study on a systematically collected 
benchmark comprising \revise{$37$} victim protocols involved in security 
hacking incidents. \revise{The scope of this work is described in 
Section~\ref{sec:preli}, and the details of our benchmark collection 
methodology are provided in Section~\ref{subsec:study_methodology}.} 
For each hacking incident, we examine the uniqueness 
of the control flows in the hack transactions, determining whether 
they differ from all 
previously observed transactions. The detailed study results are available 
in the ``RQ1'' column in Table~\ref{tab:studyTable} in 
Section~\ref{sec:evaluation}. 
In this study, we only collect the top-level function calls 
to contracts of the victim protocol, 
ignoring deeper control flows within external calls.


% \revise{
% Our hypothesis is that function-level control flow patterns can effectively 
% distinguish between benign and malicious transactions, 
% making them a practical choice for runtime security monitoring. 
% To test this hypothesis, we examine historical transactions 
% of victim DeFi protocols leading up to security incidents. 
% Typically, a victim protocol operates for a period before 
% being compromised, during which various blockchain actors 
% may build upon the protocol, introducing novel control flows. 
% Our goal is to thoroughly examine these control flows and answer 
% the following research question:
% }


% \noindent \textbf{RQ1}: 
% How do \revise{(function level)} control flows in hack transactions differ 
% from those in other 
% (benign) transactions 
% prior to a hack?





\noindent \textbf{Results:}
Out of the \revise{$37$} studied hack incidents, 
\revise{$32$} demonstrated control flows that
were distinct from any previously observed transaction patterns (marked as
\cmark in Table~\ref{tab:studyTable}), highlighting the novel mechanisms by
which these exploits were conducted. However, in $5$ cases (bZx, VisorFi, 
Opyn, DODO and
Bedrock\_DeFi), the control flows had been observed previously in benign
transactions. A detailed investigation into these exceptions revealed
insightful nuances. In particular, the Bedrock\_DeFi hack
(marked as \xmark) was traced back to an issue which mistakenly set the
conversion ratio of ETH/uniBTC to 1:1. This exploit was executed
through one function call to steal funds, which, while not novel in terms
of control flow pattern, leveraged a specific code vulnerability to 
breach the
protocol~\cite{bedrock2025}. Catching this hack involves
combining control flow analysis with data flow analysis.
The remaining $4$ hacks posed a different complexity: 
each involved various
types of re-entrancy (marked as $\halfmark$). 
Although these attacks appeared to have familiar top-level 
function calls, they triggered unique and sophisticated control 
flows at deeper interaction levels. 
This shows that, even the
simple control flow-based analysis 
can detect \revise{$32$} hacks,
it misses $4$ that require a more refined control 
flow-based approach to identify. 
Thus, we propose to enhance the control flow analysis 
framework 
to more accurately detect and classify attacks with 
intricate control flows, 
such as these $4$ re-entrancy hacks. 
The details of this enhanced 
approach are presented in Section~\ref{sec:approach}. 
Unless otherwise stated, 
we use the term ``control flow'' throughout 
the rest of this paper to refer to function-level control flow.

\noindent \textbf{Finding:} 
the vast majority of hack transactions introduce
unique function-level control flows that differ from all 
previously observed benign transactions.


